% META: DONE.1

\chapternoindex{Заключение}

В рамках данной выпускной квалификационной работы достигнута поставленная цель: разработан программный модуль «НейроКОМПАС-3D», интегрируемый со средой КОМПАС-3D и работающий локально на стороне клиента. Модуль по анализу геометрии модели позволяет реконструировать последовательность породивших её формообразующих операций с использованием методов машинного обучения.

В ходе выполнения работы решены все сформулированные задачи. Основные научные и практические результаты:

\begin{enumerate}
    \item Проведён всесторонний анализ предметной области, выявлены ключевые ограничения существующих подходов (неоднозначность интерпретации, чувствительность к шумам, отсутствие обучаемости). Обоснована необходимость гибридной архитектуры.
    
    \item Сформирован и очищен репрезентативный датасет из 10483 моделей (исходно 22275), каждая из которых содержит полную историю построения в формате JSON и геометрические представления (STEP/STL). Разработан конвейер автоматической разметки рёбер для обучения сегментационных нейросетей.
    
    \item Обоснован и реализован гибридный подход, сочетающий нейросетевую сегментацию (MeshCNN) для локализации следов операций с детерминированным вычислением параметров в CAD-ядре. Подход снимает с нейросети задачу точной регрессии параметров, чувствительной к масштабу и системе координат.
    
    \item Разработаны и обучены нейросетевые классификаторы и сегментаторы для четырёх базовых операций: \texttt{bossExtrusion} (выдавливание), \texttt{cutExtrusion} (вырезание), \texttt{fillet} (скругление), \texttt{chamfer} (фаска). На тестовых выборках достигнута точность классификации до 88.2\% (для \texttt{chamfer}) и точность сегментации (IoU) до 0,84 (для \texttt{chamfer}), доля моделей с точностью сегментации не менее 90\% составила 88,3\% для фаски и 83,7\% для скругления.
    
    \item Создан исполняемый модуль, реализующий полный цикл инференса: адаптивную триангуляцию STEP-файла, нормализацию сетки до фиксированного числа рёбер (2250), вызов обученных классификаторов/сегментаторов в автоматическом или ручном режиме, постобработку (фильтрация по площади, выделение связного компонента) и подготовку результатов (STL выделенной геометрии, JSON с типом операции и уверенностью). Модуль интегрирован с КОМПАС-3D через API и позволяет итеративно упрощать модель путём удаления распознанных операций.
\end{enumerate}

Разработанный прототип перешёл в стадию опытной эксплуатации в рамках договора с ООО «АСКОН-Системы проектирования» и показал работоспособность на наборе реальных деталей (примеры приведены в приложении Г). Модуль применим для упрощения базовых 3D-моделей, восстановления параметрической истории импортированных моделей, а также для обучения инженеров-конструкторов.

\textbf{Перспективы дальнейшего развития:}
\begin{itemize}
    \item увеличение датасета за счёт использования всех промежуточных шагов построения (потенциальное увеличение в 5-10 раз);
    \item поддержка новых типов операций (вращение, линейные и круговые массивы, кинематические операции);
    \item улучшение постобработки (фильтрация по длине рёбер, по расстоянию от центра масс);
    \item раздельное обучение для прямолинейных и криволинейных контуров с целью устранения дисбаланса плотности триангуляции;
    \item автоматическая оценка качества через API КОМПАС-3D (процент успешно упрощённых моделей, среднее количество восстановленных шагов);
    \item увеличение разрешения сетки (до 4500-6000 рёбер) и ёмкости нейросетевой архитектуры для распознавания мелких деталей.
\end{itemize}

Таким образом, представленная работа вносит вклад в решение актуальной задачи восстановления истории построения параметрических моделей и демонстрирует эффективность гибридного подхода, сочетающего машинное обучение с детерминированными алгоритмами CAD-систем. Полученные результаты могут служить основой для дальнейших исследований и промышленного внедрения.

